% Nouveau EGA V (cote 25): archive cote 25, site page 30 (PDF page 31).
% typescript. Draft compared: none.
% AI-transcribed by Claude Opus 5.5 (claude-opus-5-5, adaptive thinking, effort high), Message Batches API, batch ega25b (msgbatch_01A7n5YuXubUPNK1zsepi8RG), 2026-11-13; not yet checked by eye.
\scanpage{25/30}
\mspage{26}

\deleted{tu de III 7 la factorisation de Stein de \xout{\unsure{X}} Y $\longrightarrow$ S se fait par un}
\deleted{morphisme fini \uline{étale} S'$\longrightarrow$S, ce qui prouve que le nombre de composantes}
\deleted{connexes géom de \xout{\unsure{s}} la fibre de Y\hand{$_{\xi}$} en \xout{\unsure{t}} (égale au nombre analogue pour}\margin{Il y a une variante de \unsure{12.2.1. (4) ou (5)} \unsure{a} qui permet d'inclure ces propriétés dans 7.3.}
\deleted{S' sur S) est constant au voisinage de s, donc en l'occurence égal à 1.}
\deleted{Ce fait mérite d'être signalé au moins en remarque.}

\uline{Proposition} 7.4. Soit \unsure{$g$|$f$}: X $\longrightarrow$ S un morphisme propre et plat, \xout{et}
soit s $\in$ S, \xout{\unsure{tellxque}} posons n = dim X$_{s}$ et supposons que la dimension de
l'ènsemble des points de X$_{s}$ en lesquels X$_{s}$ n'est pas séparable sur k(s)
soit \hand{$<$} n\inserted{(\unsure{p.ex} $X_s$ séparable)}. Alors il existe un voisinage ouvert U de s, et un morphisme
étale fini et surjectif S'$\longrightarrow$U, tel que X$\mathrm{x}_{S}$S' admette une section sur S'.
Lorsque k(s) est infini, on peut prendre pour S' un sous-préschéma fermé
de X U.

Supposons d'abord k(s) infini. Nous \xout{\unsure{appliquons}} procédons par récu%
% NOTE: the end of this line and of several following lines is cut by the right edge of the scan
rence sur n, le cas n=0 étant trivial: en effet, dans ce cas il existe
un voisinage ouvert U de s tel que X$|$U lui-même soit étale, fini, et
surjectif au dessus de U, comme on voit par des cross-références immédia\xout{\unsure{t}}
tes. Lorsque n$>$0, on applique 7.3. pour la propriété ``séparable'', ce qui
nous permet de \xout{xxxxxxxxx} remplacer X par une ``section hyperplane'' Y, ayan
les mêmes propriétés, à cela près que n est rèmplacé par n-1. Lorsque
k(s) n'est pas supposé infini, on commence par faire une extension étale
de la base, ça marche.

\uline{Remarque\hand{s}} 7.5. En particulier, si X est projectif et séparable sur S,
il admet localement sur S des multisections étales. Mais on notera qu'on
peut donner des exemples, avec X propre et lisse (mais non projectif) su
S, où la même conclusion est en défaut. Bien entendu, l'hypothèse de
\xout{qua} projectivité ne peut -être affaiblie \inserted{en général} en une hypothèse quasi-projecti
comme on voit par exemple en prenant X étale non fini sur S ....

\margin{Possibilité d'éviter un ensemble \unsure{fini} ...}

\mspage{30}
